\section{Introduction}\label{sec:intro}

A structural law, in the sense this series gives the term, is a theorem about declared abstract
data --- a carrier, a quotient, a ledger, a family of patches --- so that any system supplying the
data inherits the conclusion. \FoundIV{} collected fifty-two such laws
\citep{Tsiokos2026FoundIV}. Every one of them is about fixed data: one map, one pair of quotients,
one declared finite family of patches. None quantifies over the steps of a run. Yet many of the
questions that make dynamics hard are questions about runs. Does every orbit of a map eventually
descend? Does a run terminate when the way its states are written down changes at every step?
Can a family of targets be covered when no fixed finite package of patches suffices? What happens
when the obstruction to progress is produced by the run's own history, or chosen by an adversary?
When do route choices integrate into a conserved ledger, when do local defects turn into motion,
what survives repeated loss, and how do local rules force global structure?

This paper states thirteen candidate laws, G1--G13, one for each of these questions, and sorts
them by what has actually been proved. Each entry has a general statement whose \emph{hard
hypothesis} --- a coverage property, a closed family of states, a hierarchy of forced scales ---
carries the difficult content. Seven entries are \emph{laws}: for each, the hard hypothesis is
proved here for a concrete system, although for G13 that system is only a small arithmetic example.
Two are \emph{general theorems}, whose hypotheses are ordinary data and whose classical instances
are imported. Four are \emph{templates}: the implication is proved, but its hard hypothesis is not
established in this paper for any system of interest and remains an open obligation. \Cref{tab:catalog} records the sorting, and
separately what is checked in Lean.

\subsection{A worked example: residues do not predict the Collatz ledger}\label{sec:intro-example}

The accelerated Collatz map sends an odd number $n$ to $T(n) = (3n+1)/2^{a(n)}$, where $a(n)$ is
the exponent of the largest power of $2$ dividing $3n+1$. Take $n = 5$ and $n = 9$. Both are
congruent to $1$ modulo $4$, so they agree in their last two binary digits. But $3\cdot 5 + 1 =
16$ gives $a(5) = 4$, while $3 \cdot 9 + 1 = 28$ gives $a(9) = 2$. Knowing $n$ modulo $4$ does not
determine even the first division exponent.

The same happens at every depth: for each $m \ge 1$ there are two odd numbers that agree modulo
$2^m$ whose first division exponents are $m$ and at least $m+1$ (\Cref{thm:split-pair}). The pair
above is the case $m = 2$. At the same time the bookkeeping of an orbit is exact. After $k$ steps,
with $A_k$ the sum of the first $k$ exponents,
\[
  2^{A_k}\, n_k \;=\; 3^k\, n + B_k ,
\]
for an explicit integer $B_k$ determined by the exponents (\Cref{thm:affine-ledger}). The orbit has
dropped below its start at step $k$ exactly when $2^{A_k} n > 3^k n + B_k$, an inequality between
integers. Whether every orbit eventually satisfies it is the Collatz conjecture, which this paper
does not settle. The two facts show where the difficulty lives: in an exact ledger whose next entry
is not predicted by any fixed number of low-order digits.

\subsection{Obstructions at every depth}\label{sec:intro-contribution}

Two notions recur throughout the paper. A \emph{quotient} $q_m$ records what an observer sees at
depth $m$, for example the last $m$ binary digits of a number, and a \emph{readout} $r$ records a
future fact, for example the next division exponent. A \emph{split pair} for $q_m$ and $r$ is a pair
of states with the same quotient and different readouts. If a split pair exists at every depth, no
observer who looks at a fixed number of digits can predict the readout.

The Collatz pairs above are such an obstruction for the last $m$ binary digits and the first
division exponent. The orbits of $2^L - 1$ give a second one for the time of first descent: they lie
in the residue class of $-1$ modulo $2^m$ once $L \ge m$, and they stay above their start for at
least $L - 1$ steps (\Cref{rem:collatz-paper}). So no bound on the time of first descent can depend
only on the residue of $n$ modulo a fixed power of $2$.

Each entry of the catalog is stated in one form: the data of a run, a certificate in the run's own
terms, and the named hypothesis that makes the certificate available. When that hypothesis is not
yet proved for any system of interest, the entry is a \emph{template}, and its value is that it
names the obligation precisely and lets it be recognized in unrelated systems. The split-pair
condition above, for instance, is exactly what must be avoided by any proof that a fixed-depth
reading decides a run; the same condition reappears for binary reverse-and-add, where a different
tool, a family of states closed under the map and free of palindromes, settles one orbit
(Result~B below).

\subsection{Main results}\label{sec:intro-results}

The following four results are proved in the text and checked in Lean. Their relation to the
catalog differs.
\begin{itemize}
  \item Result~A (\Cref{sec:G1}) gives the certificate that G1 asks for, in exact integer form, and
    shows that no fixed number of binary digits predicts it. It does not establish the hypotheses
    of G1.
  \item Result~B (\Cref{sec:G5}) establishes the hard hypothesis of G5 for one orbit.
  \item Result~C (\Cref{sec:G10}) is a collapse example for G10. It shows that parity cannot be the
    protected invariant for vectors of length $2^k$.
  \item Result~D (\Cref{sec:G8}) is the general statement of G8, with two examples that separate
    it from weaker properties.
\end{itemize}

\begin{headline}[A: the Collatz ledger and its split pairs]
Every orbit of the accelerated Collatz map satisfies $2^{A_k} n_k = 3^k n_0 + B_k$ for every
$k$. For every depth $m \ge 1$ there are odd positive $n_1 \equiv n_2 \pmod{2^m}$ whose first
division exponents are at least $m+1$ and exactly $m$. Hence no fixed residue class modulo $2^m$
determines the ledger (\Cref{thm:affine-ledger,thm:split-pair}).
\end{headline}

\begin{headline}[B: reverse-and-add in base 2]
Under binary reverse-and-add, the orbit of $22 = 10110_2$ reaches a four-word family after four
steps, cycles through its four phases with a parameter that grows by one per cycle, and never
reaches a palindrome (\Cref{thm:four-phase,cor:orbit-22}).
\end{headline}

\begin{headline}[C: Ducci nilpotency]
On cyclic binary vectors of length $2^k$, the map $x_i \mapsto x_i + x_{i+1}$ over
$\mathbb{F}_2$ vanishes after $2^k$ steps (\Cref{thm:ducci}).
\end{headline}

\begin{headline}[D: abelian move systems]
Consider a system of local moves in which any two distinct moves that are legal together remain
legal after each other and commute. Start from a configuration that admits no infinite legal run.
Then any two legal runs that end in stable configurations end in the same configuration, and they
fire every site the same number of times (\Cref{thm:G8}). Two small examples show that a unique
final configuration does not give canonical firing counts, and that termination does not give a
unique final configuration (\Cref{ex:G8-b,ex:G8-d}).
\end{headline}

\subsection{What is imported, and what is not claimed}\label{sec:intro-claims}

The deep theorems the catalog relies on are imported with attribution: Goodstein's theorem and the Kirby--Paris independence result, Dhar's abelian
property and the least-action principle of Fey, Levine and Peres, the solution of the angel
problem, the aperiodic tilings of Berger and Robinson, the de Bruijn--Erd\H{o}s compactness
theorem, and the density theorems of Bourgain and Kontorovich together with the reciprocity
obstruction of Haag, Kertzer, Rickards and Stange. The Collatz conjecture, the Erd\H{o}s--Straus
conjecture, the base-10 Lychrel problem, Recam\'an's coverage question, the Langton's-ant highway
conjecture, Gilbreath's conjecture and the chromatic number of the plane all remain open.

\paragraph{Roadmap.} \Cref{sec:related} places the paper in the literature, and \Cref{sec:setting}
fixes the vocabulary and gives the catalog. \Cref{sec:laws} treats the laws, \Cref{sec:theorems}
the general theorems and \Cref{sec:templates} the templates; the four headline results appear in
their entries. \Cref{sec:scope} lists inferences the results do not support,
and \Cref{sec:lean} records what is formalized. The supplement holds the source concordance, the
full register of non-implications, the case analyses and instance lists, the finite computations,
the reading of three laws as asymptotic statuses, and the replay commands.
